
\textbf{Data sources.} We extracted data from the following primary sources:
\begin{itemize}
\item GPT-4 Technical Report (OpenAI, March 2023)
\item GPT-3.5-turbo documentation (OpenAI, 2022)
\item Claude-3 Model Card (Anthropic, March 2024)
\item Claude-2 technical documentation (Anthropic, 2023)
\item Llama-3 Research Paper (Meta, April 2024)
\item Llama-2 paper (Meta, 2023)
\end{itemize}

\textbf{Implementation.} We implemented the extraction and validation pipeline in Python with the following components:
\begin{itemize}
\item Report downloading with retry logic and caching
\item Manual extraction of category-level tables with schema normalization
\item Validation logic to check coverage, granularity, and completeness against thresholds
\item Generation of summary statistics and visualizations
\end{itemize}

Execution time was approximately 5 seconds for loading curated data. No GPU resources were required. Extracted data were stored in CSV format with accompanying metadata documenting extraction decisions.

\textbf{Evaluation metrics.} We computed:
\begin{itemize}
\item Family coverage count: Number of families (out of 3) with both baseline and current data
\item Category granularity: Number of categories per benchmark
\item Data completeness: Percentage of expected data cells present, calculated as (cells present / cells expected) × 100\%, where cells expected = families × timepoints × benchmarks × categories
\item Temporal coverage: Binary indicator of whether both timepoints are present per family
\end{itemize}

These metrics are deterministic properties of the extracted data. No statistical significance testing is required, as we measure data availability rather than estimating population parameters.
